% yoshi_revisedv2.tex --- MIプログラム 中間発表予稿（A4・2段組・2ページ）
% コンパイル: platex yoshi_revisedv2.tex ; dvipdfmx yoshi_revisedv2.dvi
\documentclass[a4paper,twocolumn,10pt]{jsarticle}
\usepackage{amsmath,amssymb}
\usepackage[dvipdfmx]{graphicx}
\usepackage{booktabs}
\usepackage{url}
\setlength{\columnsep}{7mm}
\setlength{\textheight}{272mm}
\setlength{\topmargin}{-21mm}
\setlength{\textwidth}{175mm}
\setlength{\oddsidemargin}{-7mm}
\setlength{\evensidemargin}{-7mm}
\setlength{\parindent}{1zw}
\pagestyle{empty}
\setlength{\textfloatsep}{2pt plus 1pt minus 1pt}
\setlength{\floatsep}{2pt plus 1pt minus 1pt}
\setlength{\intextsep}{2pt plus 1pt minus 1pt}
\setlength{\abovecaptionskip}{1pt}
\setlength{\belowcaptionskip}{0pt}
\setlength{\abovedisplayskip}{1.5pt plus 1pt minus 1pt}
\setlength{\belowdisplayskip}{1.5pt plus 1pt minus 1pt}
\setlength{\abovedisplayshortskip}{1pt plus 1pt minus 1pt}
\setlength{\belowdisplayshortskip}{1pt plus 1pt minus 1pt}

\newcommand{\CVaR}{\mathrm{CVaR}}

\begin{document}

\twocolumn[{%
\begin{center}
{\normalsize 【発表日程】\quad MIプログラム\quad 中間発表}\par\vspace{1mm}
{\LARGE\bfseries 歪み関数重み付き1-Wasserstein距離に基づく}\par\vspace{1mm}
{\LARGE\bfseries quantile fraction学習とCVaR推定誤差の低減}\par\vspace{2mm}
{\normalsize 発表者: 【氏名】\qquad 所属: 【研究室名】\qquad 指導教員: 【指導教員名】}
\end{center}
\vspace{1mm}
}]

\section{はじめに}
金融派生商品のヘッジ等では，損益分布の下側裾に生じる極端な損失の抑制が求められる．本稿では収益（損益）$Z$（大きいほど望ましい）に対する下側 $\CVaR_\alpha(Z) = \alpha^{-1}\int_0^\alpha F_Z^{-1}(\tau)\,d\tau$（$F_Z^{-1}$ は分位点関数）\cite{rockafellar} を単に $\CVaR_\alpha$ と呼び，これを最大化する（通常の損失型CVaRとは符号が逆）．下側 $\alpha$ 領域の平均損益に基づくリスク回避方策の学習には，下側裾の正確な分布評価が必要となる．

収益分布全体を学習する分布型強化学習（分布型RL）はCVaRの評価にも応用可能であるが，IQN \cite{iqn} や FQF \cite{fqf} などの既存手法は分布全体の一様近似を前提としており，下側裾を重点的に評価する設計にはなっていない．特に FQF は，真の分位点関数 $F^{-1}$ と近似分位点関数 $\hat F^{-1}$ の間の1-Wasserstein距離 $W_1 = \int_0^1 |F^{-1}(\tau) - \hat F^{-1}(\tau)|\,d\tau$ を最小化するように quantile fraction（累積確率軸 $[0, 1]$ の区切り位置 $\tau_i$）を学習する．しかしこれは $[0,1]$ 上の一様測度に基づくため，累積確率の数％程度の下側裾 $[0, \alpha]$ に配分される重みが小さく，原子（確率質量を置く離散点）の多くが中央から上側に配分されて下側裾の解像度が不足する．例えば等間隔配置（$N=16$）では最悪5\%のリスク領域 $[0, 0.05]$ に内部の quantile fraction が1点も入らず，CVaR領域を粗い代表点でしか表現できないためCVaR推定誤差が増大する．そこで本研究では，歪み関数を目的関数に組み込み，quantile fraction の配置自体をCVaR推定に対して最適化する．

\section{関連技術・先行研究}
分位点回帰型の分布型RLは，収益分布の分位点関数を分位点回帰損失によって学習する．なかでも本研究の基礎となるFQF \cite{fqf} は図\ref{fig:fqf}に示すように，状態を入力としてquantile fraction $\tau_i$ を提案するFraction Proposal Network（FPN）と，分位点値を出力する分位点値ネットワークからなる．従来のFPNは一様重み $W_1$ を最小化するように学習される．

\begin{figure}[t]
\centering
\vspace{-2pt}
\includegraphics[width=0.82\columnwidth]{fig1_fqf_overview.png}
\vspace{-2pt}
\caption{FQFの概要．FPNがquantile fraction $\tau_i$を提案（本研究では目的関数を歪み重み付き$W_1$へ置換）．}
\label{fig:fqf}
\vspace{-2pt}
\end{figure}

歪み関数 $g$ を用いた歪みリスク尺度は，分位点関数の重み付き積分 $\rho_g(Z) = \int_0^1 F^{-1}(\tau)\,dg(\tau)$ で表され，$\CVaR_\alpha$ は $g_\alpha(\tau)=\min(\tau/\alpha,1)$ の場合に対応する．確率分布を有限個の代表点で近似する最適量子化理論（Graf and Luschgy \cite{graf}）によれば，重みに応じて原子を最適配置することで推定精度の向上が見込める．

しかし既存手法では，IQN \cite{iqn} は行動選択時のサンプリング歪みのみで原子配置を学習せず，Ganらの手法（FQF-SRM）\cite{fqfsrm} や Beyond CVaR \cite{beyondcvar} もリスク尺度を行動選択側に用いる一方，FPN損失自体を歪み重み付き1-Wasserstein距離へ置換して配置最適化するものではない．調査した限り，歪みリスク尺度に応じて quantile fraction配置自体を学習する定式化は確認できていない．

\section{研究目的}
本研究の目的は，FPNの目的関数を一様重みの $W_1$ から歪み関数重み付き $W_1$ へと置換し，歪み分布 $\mathcal{D}_g Z$（累積分布を $g$ で変換した分布）の1次元最適量子化理論に基づき quantile fraction を学習する手法を確立することである．歪みリスク尺度の推定誤差は既知のLipschitz評価 \cite{pichler2013} により上界 $L_g := W_1(\mathcal{D}_g Z,\mathcal{D}_g\hat Z)$ で抑えられる．
\begin{equation}
\bigl|\rho_g(Z)-\rho_g(\hat Z)\bigr| \le \int_0^1 \bigl|F^{-1}(\tau)-\hat F^{-1}(\tau)\bigr|g'(\tau)\,d\tau = L_g
\label{eq:bound}
\end{equation}
この $L_g$ を直接最小化することで，歪み重みに整合したCVaR推定誤差の低減を目指す．

\section{提案手法（アルゴリズム）}
目的関数 $L_g$ を最小化するため，quantile fraction $\tau_i$ に関する勾配の閉形式を導出した．各区間 $[\tau_i, \tau_{i+1}]$ の最適代表分位点は，歪み測度を二等分する $g$-中央値 $\hat\tau_i$（$g(\hat\tau_i)=\{g(\tau_i)+g(\tau_{i+1})\}/2$）での $F^{-1}(\hat\tau_i)$ で与えられる．これより $L_g$ の $\tau_i$ に関する勾配は
\begin{equation}
\frac{\partial L_g}{\partial\tau_i} = g'(\tau_i)\bigl[2F^{-1}(\tau_i)-F^{-1}(\hat\tau_i)-F^{-1}(\hat\tau_{i-1})\bigr]
\end{equation}
となる．これは従来のFQFの勾配式に $g'(\tau_i)$ を乗じた閉形式である．RL実装では各 $F^{-1}$ を分位点値ネットワークの出力で評価し，定数扱いとして quantile fraction のみを更新する形で安定的に学習できる．
なお $\tau>\alpha$ での勾配消失を防ぐため，正則化歪み $g_{\alpha,\lambda}(\tau)=(1-\lambda)g_\alpha(\tau)+\lambda\tau$（$\lambda=0.05$）で更新した（評価は純粋な $g_\alpha$ で行う．純粋 $\CVaR$ との誤差上界の差は係数約1.05倍）．

\section{現在までの進展}
\vspace{-2pt}

\subsection{実験A\quad 真値参照の1次元最適量子化}
\vspace{-1pt}
真の分位点関数が既知の設定で，配置方針（一様配置M0，$g$-一様配置M1，提案の歪み最適配置M2）によるCVaR推定誤差を比較した．対象分布は，ガウス混合 $0.9\mathcal{N}(0,1)+0.1\mathcal{N}(-4,1)$（D1），負の対数正規（D2），自由度 $\nu=3$ の $t$分布（D3）の3分布（$N\in\{8, 16, 32, 64, 128, 256\}$，$\alpha\in\{0.01, 0.05, 0.1\}$）である．

\begin{figure}[htbp]
\centering
\vspace{1pt}
\includegraphics[width=\columnwidth]{fig2_placement_crop.png}
\vspace{1pt}
\caption{D1（ガウス混合，$N=16,\alpha=0.05$）における配置比較．縦線は quantile fraction $\tau_i$，丸は代表点 $\hat\tau_i$（右は対応する収益分位点値），網掛けは下側裾 $[0,\alpha]$．M0は一様配置，M1は裾内等分割配置，M2は提案の歪み最適配置を示す．}
\label{fig:placement}
\vspace{-1pt}
\end{figure}

図\ref{fig:placement}に示すように，提案法M2では15点中13点が下側裾に集中し，代表設定（$\alpha=0.05, N=16$）における推定誤差改善比 $R = e_{\mathrm{M0}}/e_{\mathrm{M2}}$ は3分布すべてで $R\ge 39.6$（最大97.9）となった．例えばD1ではM2の16点と同等の精度を一様配置で得るには約695点を要し（D2，D3は探索上限1024点でも未達），高い点数効率を示した．

収束次数においても，一様配置M0は両対数プロットの傾きが $-0.72$〜$-1.27$ と一次精度 $O(N^{-1})$ 程度にとどまり，ヘビーテイルのD3では端点特異性（$\tau\to0$ での $F^{-1}$ の発散）により $-0.72$ まで悪化した．一方，提案法M2は端点特異性を抑えることで，傾き $-1.79$〜$-1.84$ と二次精度 $O(N^{-2})$ に近い挙動を示した．

\vspace{-2pt}
\subsection{実験B\quad 強化学習環境における検証}
\vspace{-1pt}
安全なSAFE行動と，稀に極端な損失（$t$分布）が生じる高期待値RISKY行動を選ぶ有限ホライズン（$H=10$）環境RiskyChain-v0（55状態$\times$2行動，$\gamma=1$）を構築し，$\CVaR_{0.05}$ を最大化するエージェントを100kステップ・3シードで学習した．

表\ref{tab:expB}に結果を示す．従来FQF（B0）は一様重み $W_1$ で学習するが，下側裾内の内部 quantile fraction が0点となり高リスク方策（経験CVaR $-10.79$）に収束した．目的関数のみ置換したB1（標準初期化）では，初期出力の一様性で裾外（$\tau>0.05$）の勾配重み $g'(\tau_i)$ が微弱となり，一部シードで局所解に陥った．これに対し初期バイアスを調整し初期配置を下側裾へ寄せた提案法B1g（裾初期化）では，勾配が有効に作用し下側裾へ quantile fraction が集中配分（15点中平均14.06点）された．推定MAEは $7.02\to0.96$（約86\%低減）となり，リスク回避方策（経験CVaR $0.70$）を獲得した．この差には各手法の獲得方策に伴う分布難易度の交絡も含まれるが，全時刻でSAFEを選択する参照方策（all-SAFE）下でも $1.88\to0.59$（約69\%低減）となった．

\begin{table}[t]
\centering
\vspace{-2pt}
\setlength{\tabcolsep}{2.5pt}
\caption{実験Bの性能比較（3シード平均$\pm$SD）．推定MAEは獲得方策下の$\CVaR_{0.05}$推定誤差，経験CVaRは5,000エピソードの実収益．裾点数は15個の内部 quantile fraction中$[0,\alpha]$に入る数（IQNは連続サンプリングのため---）．}
\label{tab:expB}
\resizebox{\columnwidth}{!}{%
\begin{tabular}{lccc}
\toprule
手法 & 推定MAE & 経験CVaR & 裾点数\\
\midrule
B0 従来FQF & $7.02\pm0.14$ & $-10.79\pm0.00$ & $0.00\pm0.00$\\
B1 歪みFQF（標準初期化） & $1.41\pm0.92$ & $-1.15\pm2.38$ & $0.48\pm0.46$\\
\textbf{B1g 提案法（裾初期化）} & $0.96\pm0.51$ & $0.70\pm0.43$ & $14.06\pm0.11$\\
B2 リスク感応IQN & $0.36\pm0.05$ & $1.01\pm0.01$ & ---\\
B2p IQN（歪みサンプリング） & $3.42\pm1.38$ & $1.01\pm0.02$ & ---\\
B3g 歪みFQF（裾固定） & $0.74\pm0.33$ & $0.30\pm0.83$ & $15.00\pm0.00$\\
\bottomrule
\end{tabular}%
}
\vspace{-2pt}
\end{table}

他方，暗黙的表現のIQN（B2）は，学習時は$[0,1]$全体を一様サンプリングし行動選択時のみ裾を抽出することで，本実験では最小の推定MAE（$0.36\pm0.05$）を示した．これに対し，サンプリング分布を裾に偏らせたB2p（歪みサンプリング）では，重要度補正なしのサンプリングによりベルマンターゲットが実質的に歪んだ可能性がある．これにより，方策性能（経験CVaR $1.01$）は保たれるものの推定MAEが約10倍（$3.42$）に悪化した．一方で，提案法B1gはベルマン更新（分位点回帰）と原子配置（FPN）を分離して学習するため，ベルマンターゲットの確率重みを保ちつつ下側裾の解像度を高められた．また，FPNによる適応学習を行わず下側裾へ固定配置したB3g（裾固定）も，推定MAEはSDの範囲内でB1gと同水準（$0.74\pm0.33$）となった．これは本環境では収益分布の状態依存性が小さく固定配置で十分だったためと考えられ，IQNに対する一般的優位性やFPNの状態適応性の実証は今後の課題である．

\section{今後の展望}
\vspace{-2pt}
\begin{enumerate}
\setlength{\topsep}{0pt}
\setlength{\itemsep}{-1pt}
\setlength{\parskip}{0pt}
\setlength{\parsep}{0pt}
\item \textbf{金融環境での実証（実験C）}：動的環境でFPNの状態適応性が固定配置を上回ることを実証する．
\item \textbf{収束特性の定理化}：適応格子の端点特異性緩和を理論評価し，改善比下界 $\Omega(1/\alpha)$ の証明を目指す．
\item \textbf{共通方策下の公平比較}：評価プロトコル統一により状態適応性と下側裾集中配置の寄与を分離評価する．
\end{enumerate}
\vspace{-2pt}

\begin{thebibliography}{9}
\fontsize{6pt}{7.2pt}\selectfont
\setlength{\itemsep}{-0.8pt}
\setlength{\parskip}{0pt}
\bibitem{rockafellar} R. T. Rockafellar, S. Uryasev. Optimization of Conditional Value-at-Risk. Journal of Risk, 2(3):21--41, 2000.
\bibitem{iqn} W. Dabney et al. Implicit Quantile Networks for Distributional Reinforcement Learning. ICML, 2018.
\bibitem{fqf} D. Yang et al. Fully Parameterized Quantile Function for Distributional Reinforcement Learning. NeurIPS, 2019.
\bibitem{graf} S. Graf, H. Luschgy. Foundations of Quantization for Probability Distributions. Springer, 2000.
\bibitem{fqfsrm} G. Gan et al. Energy-aware job scheduling for green data centers under uncertainty. Applied Energy, 420:128137, 2026.
\bibitem{beyondcvar} M. Moghimi, H. Ku. Beyond CVaR: Exploring General Static Spectral Risk Measures for Enhanced Decision-Making in Distributional Reinforcement Learning. ICML, 2025.
\bibitem{pichler2013} A. Pichler. Evaluations of Risk Measures for Different Probability Measures. SIAM J. Optim., 23(1):530--551, 2013.
\end{thebibliography}

\end{document}
